
\subsection{哈密顿同痕及双模同态的复合}
\label{subsec:5-3-4}

本小节证明定理 5.3.39，
并用它完成定理 5.3.14 的证明。
定理 5.3.27 已表明，
在双模同态的同伦意义下，
$\phi$ 不依赖
$\{J_{\tau,t}\}_{\tau,t}$、
多值截面 $\mathfrak s$
及 $\mathcal X_3(R'_{h'})$
的选择。
因此，以下将
$\phi^{\psi^{(0)},\psi^{(1)},
\{J_{\tau,t}\}_{\tau,t},\mathfrak s}$
简写为
$\phi^{\psi^{(0)},\psi^{(1)}}$。

令
$\{\psi^{(0),(j)}_\rho\}_{0\le\rho\le1}$、
$\{\psi^{(1),(j)}_\rho\}_{0\le\rho\le1}$
（$j=I,II$）
为哈密顿同痕，
且
$\psi^{(0),(j)}_0
=\psi^{(1),(j)}_0=\mathrm{id}$。
对 $i=0,1$，定义复合
$\psi^{(i),(II)}\#\psi^{(i),(I)}$ 为
\[
\bigl(
\psi^{(i),(II)}
\#\psi^{(i),(I)}
\bigr)_\rho
=
\begin{cases}
\psi^{(i),(I)}_{2\rho},
&\rho<1/2,\\
\psi^{(i),(II)}_{2\rho-1}
\circ\psi^{(i),(I)}_1,
&\rho\ge1/2.
\end{cases}
\]

\begin{samepage}
\begin{theorem}{5.3.39}\label{thm:5-3-39}
双模同态
\[
\phi^{
\psi^{(0),(II)}\#\psi^{(0),(I)},
\psi^{(1),(II)}\#\psi^{(1),(I)}}
\]
与复合
\[
\phi^{\psi^{(0),(II)},\psi^{(1),(II)}}
\circ
\phi^{\psi^{(0),(I)},\psi^{(1),(I)}}
\]
同伦。
把 $\phi$ 换成 $\varphi$，
结论也成立。
\end{theorem}
\end{samepage}

\begin{proof}
对 $i=0,1$，置
\[
\begin{aligned}
L^{(i)\prime}
&=\psi^{(i),(I)}_1(L^{(i)}),\\
L^{(i)\prime\prime}
&=\psi^{(i),(II)}_1(L^{(i)\prime})
=\psi^{(i),(II)}_1
\circ\psi^{(i),(I)}_1(L^{(i)}).
\end{aligned}
\]
考虑两段同态
\[
\begin{aligned}
\phi^{\psi^{(0),(I)},\psi^{(1),(I)}}\colon&
\bigl(
C(L^{(1)},L^{(0)};
\Lambda_{\mathrm{nov}}),
\mathfrak n^{J_t,\mathfrak s}_{k_1,k_0}
\bigr)\\
&\longrightarrow
\bigl(
C(L^{(1)\prime},L^{(0)\prime};
\Lambda_{\mathrm{nov}}),
\mathfrak n^{J'_t,\mathfrak s'}_{k_1,k_0}
\bigr),\\[1mm]
\phi^{\psi^{(0),(II)},\psi^{(1),(II)}}\colon&
\bigl(
C(L^{(1)\prime},L^{(0)\prime};
\Lambda_{\mathrm{nov}}),
\mathfrak n^{J'_t,\mathfrak s'}_{k_1,k_0}
\bigr)\\
&\longrightarrow
\bigl(
C(L^{(1)\prime\prime},L^{(0)\prime\prime};
\Lambda_{\mathrm{nov}}),
\mathfrak n^{J''_t,\mathfrak s''}_{k_1,k_0}
\bigr).
\end{aligned}
\]

令 $J^j_{\tau,t}$（$j=I,II$）
分别为构造
$\phi^{\psi^{(0),(j)},\psi^{(1),(j)}}$
所用的相容殆复结构族。\translatornote{原书实体页 329
把下列第二段双模同态的源靶均按
$(L^{(0)},L^{(1)})$ 倒序书写，
与定理所需的复合不相容。
这里依第一段的靶及既定双模次序
改为 $(L^{(1)},L^{(0)})$；
见疑误 A-0302。}
它们满足
\[
\begin{cases}
J^I_{\tau,t}=J_t,
&\tau<-A,\\
J^I_{\tau,t}=J'_t,
&\tau>A,\\
J^{II}_{\tau,t}=J'_t,
&\tau<-A,\\
J^{II}_{\tau,t}=J''_t,
&\tau>A.
\end{cases}
\]
取 $B>100A$，定义
\[
J^B_{\tau,t}
=
\begin{cases}
J^I_{\tau+B,t},&\tau\le0,\\
J^{II}_{\tau-B,t},&\tau\ge0.
\end{cases}
\]
对任意这样的 $B$，
均可用 $J^B_{\tau,t}$
定义
$\phi^{
\psi^{(0),(II)}\#\psi^{(0),(I)},
\psi^{(1),(II)}\#\psi^{(1),(I)}}$；
由定理 5.3.27，
所得 $A_\infty$ 双模同态
在同伦意义下不依赖 $J^B_{\tau,t}$。

胶合论证给出如下模空间同构。
\translatornote{原书实体页 329
把下一式中间端的纤维积底空间写成
$R''_{h''}$；
两侧评价映射都落在第一段形变后的
$(L^{(1)\prime},L^{(0)\prime})$
的相交空间，故改为
$R'_{h''}$；见疑误 A-0304。}
\[
\begin{aligned}
&\mathcal M^\circ_{k_1,k_0}
\bigl(
L^{(1)},L^{(0)};
\{J^B_{\tau,t}\}_{\tau,t};
[h,w],[h',w'];
\operatorname{top}(\tau)
\bigr)\\
&\quad\cong
\coprod_{[h'',w'']}
\coprod_{\substack{
k_1^I+k_1^{II}=k_1\\
k_0^I+k_0^{II}=k_0}}
\Bigl(
\mathcal M^\circ_{k_1^{II},k_0^{II}}
\bigl(
L^{(1)\prime},L^{(0)\prime};
\{J^{II}_{\tau,t}\}_{\tau,t};
[h'',w''],[h',w'];
\operatorname{top}(\tau)
\bigr)\\
&\qquad\qquad
\times_{R'_{h''}}
\mathcal M^\circ_{k_1^I,k_0^I}
\bigl(
L^{(1)},L^{(0)};
\{J^I_{\tau,t}\}_{\tau,t};
[h,w],[h'',w''];
\operatorname{top}(\tau)
\bigr)
\Bigr).
\end{aligned}
\]
这个同构也能延拓到紧致化的模空间。
因此，可对模空间
\[
\mathcal M_{k_1,k_0}
\bigl(
L^{(1)},L^{(0)};
\{J^B_{\tau,t}\}_{\tau,t};
[h,w],[h',w']
\bigr)
\]
选取与该同构相容的多值截面。

其紧致化见图 5.3.11。
例如，该图对应下列项。
\translatornote{原书实体页 330
的这一示例项及图中右段标记，
上、下侧的 $(0)$、$(1)$
与前文逐侧作用的记号有疑点。
目前保留原样，不擅自重排；
见疑误 A-0303。}
\[
\begin{aligned}
\phi^{II}_{3,2}\Bigl(
&\psi^{(0)}_{1*}(P_1^{(1)}),
\psi^{(0)}_{1*}(P_2^{(1)}),
\mathfrak f_1^{\psi^{(0)}}
(P_3^{(1)}),\\
&\phi^I_{2,2}
(P_4^{(1)},P_5^{(1)},S,
P_1^{(0)},P_2^{(0)}),
\psi^{(1)}_{1*}(P_3^{(0)}),
\mathfrak f_0^{\psi^{(1)}}(1)
\Bigr).
\end{aligned}
\]
\begin{figure}[htbp]
\centering
\setlength{\unitlength}{0.85mm}
\begin{picture}(145,94)(0,-14)
 \qbezier(5,56)(12,47)(20,40)
 \qbezier(20,40)(46,5)(75,40)
 \qbezier(75,40)(101,73)(130,40)
 \qbezier(130,40)(138,31)(145,26)
 \qbezier(5,24)(12,33)(20,40)
 \qbezier(20,40)(46,74)(75,40)
 \qbezier(75,40)(102,6)(130,40)
 \qbezier(130,40)(137,49)(145,55)
 \put(45,66){\circle{18}}
 \put(36,16){\circle{18}}
 \put(107,66){\circle{18}}
 \put(101,15){\circle{20}}
 \put(83,5){\circle{12}}
 \put(118,4){\circle{12}}
 \put(45,75){\makebox(0,0){$\times$}}
 \put(56,56){\makebox(0,0){$\times$}}
 \put(55,27){\makebox(0,0){$\times$}}
 \put(36,7){\makebox(0,0){$\times$}}
 \put(97,54){\makebox(0,0){$\times$}}
 \put(121,31){\makebox(0,0){$\times$}}
 \put(83,-1){\makebox(0,0){$\times$}}
 \put(118,-2){\makebox(0,0){$\times$}}
 \put(107,66){\makebox(0,0){$\scriptscriptstyle\rho\le1/2$}}
 \put(101,15){\makebox(0,0){$\scriptscriptstyle\rho\ge1/2$}}
 \put(83,5){\makebox(0,0){$\scriptscriptstyle\rho<1/2$}}
 \put(118,4){\makebox(0,0){$\scriptscriptstyle\rho\ge1/2$}}
 \put(1,58){$\scriptstyle L^{(0)}$}
 \put(1,18){$\scriptstyle L^{(1)}$}
 \put(135,57){$\scriptstyle L^{(1)}$}
 \put(135,20){$\scriptstyle L^{(0)}$}
 \put(18,44){$\scriptstyle h$}
 \put(72,45){$\scriptstyle h''$}
 \put(127,44){$\scriptstyle h'$}
 \put(40,77){$\scriptstyle P_5^{(1)}$}
 \put(57,58){$\scriptstyle P_4^{(1)}$}
 \put(49,19){$\scriptstyle P_2^{(0)}$}
 \put(23,1){$\scriptstyle P_1^{(0)}$}
 \put(91,58){$\scriptstyle P_3^{(0)}$}
 \put(116,23){$\scriptstyle P_1^{(1)}$}
 \put(70,-10){$\scriptstyle P_3^{(1)}$}
 \put(115,-11){$\scriptstyle P_2^{(1)}$}
\end{picture}
\par\smallskip
\renewcommand{\thefigure}{5.3.11}
\refstepcounter{figure}
{\normalfont\bfseries 图~\thefigure}\label{fig:5-3-11}
\end{figure}
于是定理 5.3.39 得证。
\end{proof}

至此已证明定理 5.3.27 和
定理 5.3.39。
现在用它们说明
$\phi^{
\psi^{(0)},\psi^{(1)},
\{J_{\tau,t}\}_{\tau,t},\mathfrak s}$
是同伦等价。
具体地，
分别取
$\psi^{(0)}$、$\psi^{(1)}$
的逆同痕族
$\psi^{(0)-1}$、$\psi^{(1)-1}$，
以及相应的逆殆复结构族
$\{J^{-1}_{\tau,t}\}_{\tau,t}$
和逆多值截面
$\mathfrak s^{-1}_{\tau,t}$。
用这些数据构造
\[
\phi^{
\psi^{(0)-1},\psi^{(1)-1},
\{J^{-1}_{\tau,t}\}_{\tau,t},
\mathfrak s^{-1}}.
\]
由定理 5.3.27、5.3.39，
它是
$\phi^{
\psi^{(0)},\psi^{(1)},
\{J_{\tau,t}\}_{\tau,t},\mathfrak s}$
的同伦逆
（另见注 7.2.128）。
定理 5.3.14 的证明至此完成。
